\documentclass[10pt,a4paper]{article}
\usepackage{amsmath,amssymb,amsthm,graphicx}
\usepackage{amsmath,amssymb,amsthm,graphicx,xcolor,mathrsfs,mathtools}
\usepackage[shortlabels]{enumitem}
\setlist[enumerate,1]{label={\rm(\arabic*)},ref={\rm\arabic*}}
\usepackage[a4paper,margin=25mm]{geometry}
\usepackage[colorlinks=true,linkcolor=blue,citecolor=blue,urlcolor=blue]{hyperref}
\emergencystretch=3em
\newcommand{\doi}[1]{\href{https://doi.org/#1}{doi:\nolinkurl{#1}}}
\newcommand{\arXiv}[1]{\href{https://arxiv.org/abs/#1}{arXiv:\nolinkurl{#1}}}
\providecommand{\bysame}{\leavevmode\hbox to3em{\hrulefill}\thinspace}
\newtheorem{thm}{Theorem}[section]
\newtheorem{proposition}[thm]{Proposition}
\newtheorem{corollary}[thm]{Corollary}
\newtheorem{lem}[thm]{Lemma}
\newtheorem{conjecture}[thm]{Conjecture}

\theoremstyle{definition}
\newtheorem{defin}[thm]{Definition}

\theoremstyle{remark}
\newtheorem{remark}[thm]{Remark}
\newtheorem{setting}[thm]{Setting} % definition? 
\newtheorem{question}[thm]{Question}
\newtheorem{example}[thm]{Example}


%%%%%%%%%%%%%%%%%%%%%%%%%%%%%%%%%%%%%
% Define any new command you need
%%%%%%%%%%%%%%%%%%%%%%%%%%%%%%%%%%%%%

\newcommand{\simC}{\mathrel{\underset{\scriptscriptstyle{\mathbb{P}^1}}{\sim}}}
\newcommand{\simQ}{\mathrel{\underset{\scriptscriptstyle{\mathbb{Q}}}{\sim}}}
\newcommand{\simQC}{\mathrel{\underset{\scriptscriptstyle{\mathbb{Q},
C}}{\sim}}}
\newcommand{\EEE}{\mathscr{E}}
\newcommand{\OOO}{\mathscr{O}}
\newcommand{\HHH}{\mathscr{H}}
\newcommand{\LLL}{\mathscr{L}}
\newcommand{\PPP}{\mathscr{P}}
\newcommand{\FFF}{\mathscr{P}}

\newcommand{\xExt}{\mathscr{E}xt}
\newcommand{\xHom}{\mathscr{H}om}
 \newcommand{\x}{\mathscr}
% \newcommand{\D}{\mathbb D}
 \newcommand{\N}{\mathbb N}
 \newcommand{\PP}{\mathbb P}
 \newcommand{\A}{\mathbb A}
 \newcommand{\Q}{\mathbb Q}
 \newcommand{\R}{\mathbb R}
 \newcommand{\Z}{\mathbb Z}
  \newcommand{\C}{\mathbb C}
 \newcommand{\bir}{\dashrightarrow}
 \newcommand{\rddown}[1]{\left\lfloor{#1}\right\rfloor} % round-down
 \newcommand{\nrd}[1]{\langle{#1}\rangle}

\let\oldemptyset\emptyset
\let\emptyset\varnothing

\newcommand{\pr}{\mathrm{pr}}

\newcommand{\lra}{\longrightarrow}
\DeclareMathOperator{\Cr}{Cr}
\DeclareMathOperator{\Cl}{Cl}
\DeclareMathOperator{\Aut}{Aut}
\DeclareMathOperator{\GL}{GL}
\DeclareMathOperator{\PGL}{PGL}
\DeclareMathOperator{\Pic}{Pic}
\DeclareMathOperator{\Bir}{Bir}
\DeclareMathOperator{\Sing}{Sing}
\DeclareMathOperator{\Ker}{Ker}
\DeclareMathOperator{\Exc}{Exc}
\DeclareMathOperator{\mult}{mult}
\DeclareMathOperator{\rk}{rk}


\newcounter{foo}
\makeatletter
\newcommand{\otherlabel}[2]{\protected@edef\@currentlabel{#2}\label{#1}}
\makeatother

\newcommand{\supth}[1]{\ensuremath{#1^{\mathrm{th}}}}

\title{Cyclic kernel of a K3-type action bounded by canonical Fano Weil index (Proposition 1.9)}
\author{Konstantin Loginov}
\date{arxiv:2408.11645; CC BY-SA 4.0}
\begin{document}
\maketitle
\noindent\textit{Editorial scope.} One author proposition and full current supporting proof chain; no mathematical repair, independent certification or formal admission.\par\medskip
\noindent\textit{Allowed standard prerequisites (editorial).}\begin{itemize}
\item Faithful tangent action at a fixed point. Popov [Pop14], Lemma 4; author Lemma 2.6 (lines 491–494). Standard input to the fully reproduced fixed-divisor and cyclic-kernel proofs.
\item MMP singularities, inversion of adjunction, adjunction, vanishing and Serre duality. Kollar–Mori [KM98], especially Corollary 5.25 explicitly cited at line 2730; author lines 2725–2740. External tools, while the actual quotient construction, torsion contradiction and canonical discrepancy calculation are retained.
\item Cyclic branched-cover divisibility. J. J. Wavrik, Deformations of branched coverings of complex manifolds, Amer. J. Math. 90 (1968), 926–960, Theorem 1.2; DOI 10.2307/2373491; author citation [Wav68] at line 2742. Applied by the author to the smooth locus Y\textasciicircum{}0, giving m B\textasciicircum{}0 linearly equivalent to D\textasciicircum{}0. The passage Pic(Y\textasciicircum{}0)=Cl(Y\textasciicircum{}0)=Cl(Y) and closure B is reproduced, not supplied editorially.
\item Bounded canonical Fano index. Birkar [Bir21], DOI 10.4007/annals.2021.193.2.1; Jiang–Liu [JL25], arXiv:2508.16364v1; Suzuki [Suz04], DOI 10.1007/s00229-004-0442-4. Existence of the maximal index is standard external boundedness input; Remark 11.4 retains the author’s bounds 19 (terminal case) and 66. No effective-index theorem is separately counted.
\end{itemize}\medskip
\noindent\textbf{Source and scope.} Konstantin Loginov, \emph{Finite abelian groups acting on rationally connected threefolds I: Groups of product type}, arXiv:2408.11645v3 (22 August 2026), DOI 10.46298/epiga.2026.15051. Exact author-native source under \href{https://creativecommons.org/licenses/by-sa/4.0/}{CC BY-SA 4.0}; this extracted adaptation is distributed under the same license with attribution. Changes are selection, assembly, standard formatting and restored source counters; mathematical wording is unchanged.
\par\smallskip
\noindent Only Proposition 1.9 is the counting target. All retained preliminary results, Lemma 11.2, Proposition 11.3 and Remark 11.4 are support, not additional items. The complete current quotient and torsion-elimination arguments are reproduced. Source numbering is retained across omitted passages; the full original bibliography is included to retain author citation labels and repeated-author conventions.
\par\smallskip
\noindent This is author-material preparation, not independent mathematical certification or final global-unique admission. Later in-flight overlap remains unconfirmed. Journal styles and logos are excluded; the unmodified author TeX is retained as the exact source record, and its journal dependencies are not needed by this wrapper.
\par\medskip

\section{Introduction}
We work over the field of complex numbers $\mathbb{C}$. 

Finite abelian groups are among the most elementary objects studied in algebra. In algebraic geometry, rational varieties form a rather simple and well-behaved class of varieties. Consider the following question. 

\begin{question}
\label{main-question}
Which finite abelian groups can faithfully and biregularly act on rational (or rationally connected) varieties of a given dimension? 
\end{question}

Since we consider varieties endowed with the action of a finite group, there exists an equivariant resolution of singularities, so we can consider the class of finite abelian groups that act on rational (or rationally connected) smooth projective varieties. Also,  the notions of biregular and birational group actions would give us the same list of groups in answer to Question~\ref{main-question} since we can resolve the indeterminacy of the group action.

In this paper, by a K3 surface we mean a normal projective surface $S$ with at worst canonical (that is, du Val) singularities such that $\mathrm{H}^1(S, \OOO_S)=0$ and $K_S\sim 0$.  We say that a finite abelian group~$G$ is of \emph{K3 type} if~$G$ faithfully acts on a $G\mathbb{Q}$-Fano threefold $X$ preserving a K3 surface $S\in|-K_X|$ with at worst du Val singularities. In particular, we have that $G$ is an extension of a finite abelian group $H$ that faithfully acts on a K3 surface by a cyclic group (cf. Corollary~\ref{cor-G-preserves-K3-type}):
\[
0\lra \mathbb{Z}/m\lra G \lra H\lra 0.
\]
Note that if a finite group faithfully acts on a singular K3 surface, then this action lifts to its minimal resolution.  Our main result is the following theorem (see Section~\ref{subsec-mfs} for the definitions of a $G\mathbb{Q}$-Fano variety and a $G\mathbb{Q}$-Mori fiber space).

Our next result is a boundedness property for groups of K3 type.  By the Weil index of a Fano variety $X$ with at worst canonical singularities, we mean the number $C(X)=\max\{ k\ |\, -K_X \sim k H\}$, where $H$ is a Weil divisor on $X$.

\setcounter{thm}{8}

\begin{proposition}
\label{thm-second-main-thm}
Assume that $G$ is a group of K3 type. Then in the corresponding exact sequence\color{black}
\[
0\lra \mathbb{Z}/m\lra G \lra H\lra 0, 
\] 
where the group $H$ faithfully acts on a K3 surface, we have $m\leq C$, where $C$ is  the maximal Weil index of a  Fano threefold with at worst canonical singularities.   
\end{proposition}

\setcounter{section}{2}\setcounter{subsection}{0}

 All the varieties are projective and defined over the field of complex numbers $\mathbb{C}$, unless stated otherwise. We use the language of the minimal model program (the MMP for short); see \textit{e.g.}~\cite{KM98}.

\subsection{Contractions} By a \emph{contraction} we mean a projective
morphism $f\colon X \to Y$ of normal varieties such that~$f_*\OOO_X =
\OOO_Y$. In particular, $f$ is surjective and has connected fibers. A
\emph{fibration} is defined as a contraction $f\colon X\to Y$ such
that $\dim Y<\dim X$.

Let $G$ be a finite group. By a \emph{$G$-variety} we mean a variety
$X$ together with an action of a group $G$. If $f$ is a
$G$-equivariant contraction (resp.\ a $G$-equivariant fibration) of
$G$-varieties, we call it a \emph{$G$-contraction} (resp.\ a
\emph{$G$-fibration}).

\subsection{Pairs and singularities} 
\label{subsec-pairs-sing}
A \emph{pair} $(X, B)$ consists of a normal variety $X$ and a boundary Weil $\mathbb{Q}$-divisor $B$ with coefficients in $[0, 1]$ such that $K_X + B$ is $\mathbb{Q}$-Cartier.  Let $\phi\colon W \to X$ be a log resolution of $(X,B)$, and let
\[
K_W +B_W = \phi^*(K_X +B)
\]
be the log pullback of $(X, B)$.  The \emph{log discrepancy} of a prime divisor $D$ on $W$ with respect to $(X, B)$ is $1 - \mathrm{coeff}_D B_W$, and it is denoted by $a(D, X, B)$. We say $(X, B)$ is lc (resp.\ klt, $\epsilon$-lc) if $a(D, X, B)\geq 0$ (resp.\ $> 0$, $\geq \epsilon$) for every $D$. Note that if $(X, B)$ is $\epsilon$-lc, then automatically $\epsilon \leq 1$ because $a(D, X, B) = 1$ for almost all $D$. We say that the pair $(X, B)$ is \emph{strictly lc} if it is lc but not klt.  We say that the pair $(X,B)$ is \emph{plt} if $a(D, X, B)>0$ holds for any $\phi$-exceptional divisor $D$ and for any log resolution $\phi$.  We say that the pair $(X,B)$ is \emph{dlt} if $a(D, X, B)>0$ holds for any $\phi$-exceptional divisor $D$ and for some log resolution~$\phi$.

An \emph{lc-place} of $(X, B)$ is a prime divisor $D$ over $X$, that
is, on a birational model of $X$, such that $a(D, X, B) = 0$. An
\emph{lc-center} is the image on $X$ of an lc-place.

\subsection{Log Fano and log Calabi--Yau pairs}
\label{sect-log-CY}
Let $(X, B)$ be an lc pair. We say $(X,B)$ is a \emph{log Fano pair} if $-(K_X+B)$ is ample.  If $B=0$, then $X$ is called a \emph{Fano variety}.  We say that an lc pair $(X,B)$ is a \emph{log Calabi--Yau pair} if $K_X + B \sim_{\mathbb{Q}} 0$. If $B=0$, then $X$ is called a \emph{Calabi--Yau variety}.

\subsection{Mori fiber space}
\label{subsec-mfs}
Let $G$ be a finite group.  Recall that a $G$-variety $X$ is called $G\mathbb{Q}$-factorial if every $G$-invariant Weil divisor on~$X$ is $\mathbb{Q}$-Cartier.  A $G\mathbb{Q}$-\emph{Mori fiber space} is a $G\mathbb{Q}$-factorial variety $X$ (we assume that the action of $G$ on $X$ is faithful) with at worst terminal singularities, together with a $G$-contraction $f\colon X\to Z$ to a variety $Z$ such that $\rho^G(X/Z)=1$ and $-K_X$ is ample over $Z$. If $X$ is $G$-factorial (\textit{e.g.}~if $X$ is smooth), we call it a $G$-\emph{Mori fiber space}.  If $Z$ is a point, we say that $X$ is a $G\mathbb{Q}$-Fano variety.  If $\dim X=2$ and $\dim Z=1$, then we say that $f\colon X\to Z$ is a \emph{$G$-conic bundle}.  If $\dim X=2$ and $Z$ is a point, then we say that $f\colon X\to Z$ is a \emph{$G$-del Pezzo surface}.

\setcounter{subsection}{6}\setcounter{thm}{3}

\subsection{Group actions on algebraic varieties.}
We collect some generalities on the actions of finite groups on algebraic varieties. 

\begin{defin}
Let $G$ be a finite group. By $\mathfrak{r}(G)$, we denote the \emph{rank} of $G$, that is, the minimal number of generators of $G$.
\end{defin}

\begin{remark}
\label{rem-rank-subgroup}
Note that if $G$ is a finite abelian group and $H\subset G$ is a subgroup, then $\mathfrak{r}(H)\leq \mathfrak{r}(G)$. 
\end{remark}

\begin{lem}[\textit{cf.} {\cite[Lemma 4]{Po14}}]
\label{lem-faithful-action}
Let $X$ be an algebraic variety, and let $G$ be a finite group such that $G\subset \Aut(X)$. Assume $P\in X$ is a fixed point of\, $G$.  Then the induced action of\, $G$ on the tangent space $T_P X$ is faithful.
\end{lem}

\begin{corollary}
\label{cor-faithful-action}
If a finite abelian group $G\subset \Aut(X)$ has a fixed point $P\in X$, where $X$ is an algebraic variety of dimension $n$ and~$P$ is a smooth point, then $\mathfrak{r}(G)\leq n$.
\end{corollary}

In what follows, we will need the following results. 

\begin{lem}[\textit{cf.} {\cite[Lemma 2.6]{Pr11}}]
\label{lem-fixed-curve-surface}
Let $X$ be a $3$-dimensional algebraic variety $X$ with isolated singularities, and let $G$ be a finite abelian group such that $G\subset \Aut(X)$.
\begin{enumerate}
\item\label{lfcs-1} 
If there is a curve $C\subset X$ of\, $G$-fixed points, then $\mathfrak{r}(G) \leq 2$.
\item\label{lfcs-2} 
If there is a $($possibly, reducible$)$ divisor $S\subset X$ of\, $G$-fixed points, then $\mathfrak{r}(G)\leq 1$.  If moreover $S$ is singular along a curve, then $G$ is trivial.
\end{enumerate}
\end{lem}

\begin{proof}
We prove the first claim. Let $P$ be a general point in $C$. Since $X$ has isolated singularities, $P$ is smooth on $X$. In particular, $\dim T_{P}X=3$. By Lemma~\ref{lem-faithful-action}, the induced action of $G$ on $T_{P}X$ is faithful. On the other hand, the action of $G$ on the $1$-dimensional subspace $T_P C \subset T_{P}X$ is trivial, since by assumption the action of $G$ is trivial on $C$. Hence $G$ faithfully acts on a $G$-invariant complement $T'$ to $T_P C$ in $T_P X$. So we have $\dim T'=2$. Since $G$ is a finite abelian group, it follows that $\mathfrak{r}(G)\leq 2$ in this case.  

The first part of the second claim is proven analogously. We prove the second part.  By assumption, the action of $G$ on $S$ it trivial. It follows that the induced action of $G$ on $T_P S$ is trivial for any point $P\in S$. For a general point $P\in \Sing(S)$, we have $T_P S=T_P X$ since $X$ is smooth at $P$ and $P$ is singular on $S$. By Lemma~\ref{lem-faithful-action}, we conclude that $G$ is trivial.
\end{proof}

\begin{corollary}
\label{cor-G-preserves-K3-type}
If a finite abelian group $G$ faithfully acts on a $3$-dimensional algebraic variety~$X$ with isolated singularities preserving a surface $S$, then there is an exact sequence
\[
0\lra \mathbb{Z}/m \lra G \lra H \lra 0
\]
for $m\geq 1$, where $H\subset \Aut(S)$. 
\end{corollary}

\begin{proof}
Denote the image of a natural map $\phi\colon G\to \Aut(S)$ by $H$. By Lemma~\ref{lem-fixed-curve-surface}\eqref{lfcs-2}, the kernel of $\phi$ is cyclic. The claim follows. 
\end{proof}

\setcounter{section}{10}\setcounter{thm}{0}

\section{Groups of K3 type}
\label{sect-action-on-K3}
In this section, we work with finite abelian groups of K3 type. Recall that by a K3 surface we mean a normal projective surface
$S$ with at worst canonical (that is, $1$-lc) singularities such that
$\mathrm{H}^1(S, \OOO_S)=0$ and $K_S\sim 0$. 

\begin{defin}
We say that a finite abelian subgroup $G$ is of \emph{K3 type} if $G$ faithfully acts on a $G\mathbb{Q}$-Fano threefold $X$ preserving a K3 surface $S\in|-K_X|$ with at worst du Val singularities.
\end{defin}
\iffalse
In particular, we have that $G$ is an extension of a finite abelian group $H$ that faithfully acts on a K3 surface by a cyclic group (cf. Corollary \ref{cor-G-preserves-K3-type}):
\[
0\lra K \lra G \lra H\lra 0.
\]
where $H$ is a finite abelian group that faithfully acts on a K3 surface, and 
$K=\mathbb{Z}/m$ for some $m\geq 1$.\color{black}
  \fi
  
Let $X$ be a $G\mathbb{Q}$-Fano threefold, where $G$ is a finite abelian group and $G$ acts on $X$ faithfully. Assume that there exists a $G$-invariant element $S\in |-K_X|$ that is a K3 surface with at worst canonical singularities. Consider the  exact sequence
\[
0\lra K \lra G\lra H\lra 0, 
\]
where $H\subset \Aut(S)$ and $K$ fixes $S$ pointwise.  Note that $K=\mathbb{Z}/m$ for some $m\geq 1$ (\textit{cf.}~Corollary~\ref{cor-G-preserves-K3-type}).  In particular, the group $G$ is of K3 type.  Let $Y=X/K$ be the quotient, and let $f\colon X\to Y$ be the natural projection.

\begin{lem}
\label{lem-for-thm-2}
Assume that $K$ is non-trivial. 
Then the action of\, $K$ is free in codimension~$1$ outside $S$. 
Hence, the divisorial part of the ramification of $f$ is precisely $S$. 
\end{lem}

\begin{proof}
Let $K'$ be a non-trivial subgroup of $K$. 
Assume there exists a $K'$-fixed divisor $D$ such that $D\neq S$. 
Then $K'$ fixes $S\cup D$ pointwise. 
Since $S$ is ample, we have $S\cap D\neq\emptyset$. Since $S$ is $\mathbb{Q}$-Cartier, there exists a curve $C\subset S\cap D$. At a general point of $C$, the variety $X$ is smooth, so the tangent space to $S\cup D$ at some point $P\in C$ is equal to the tangent space to $X$. However, this contradicts Lemma~\ref{lem-faithful-action} applied to the group~$K'$.
\end{proof}

\begin{proposition}
\label{prop-K3-boundedness}
The threefold $Y=X/K$ is an $H\mathbb{Q}$-factorial Fano variety with at worst canonical singularities. Moreover, $-K_Y\sim mB$ for some $B\in \Cl(Y)$.
\end{proposition}

\begin{proof}
Recall that we have $K=\mathbb{Z}/m$ for some $m\geq 1$. We may assume that $m>1$. 
First we show that $Y$ is $H\mathbb{Q}$-factorial. Indeed, this follows from the fact that
\[
\Cl(Y)\otimes\mathbb{Q} = (\Cl(X)\otimes\mathbb{Q})^{K}
\]
and hence 
\[
(\Cl(Y)\otimes\mathbb{Q})^{H} = (\Cl(X)\otimes\mathbb{Q})^{G} = (\Pic(X)\otimes\mathbb{Q})^{G}.
\]

Put $D=f(S)$. Since $K$ fixes $S$ pointwise, we have $S\simeq D$. The Hurwitz formula yields
\[
0\sim K_X + S = f^*( K_Y + D ),
\]
and hence $m(K_Y + D)\sim 0$. Note that $Y$ is a Fano variety with klt singularities. By inversion of adjunction, the pair $(Y,D)$ is plt. Observe that the divisor $D$ is $\mathbb{Q}$-Cartier.
Put $T=K_Y+D$. We show that $T\sim0$. Assume otherwise. Since $mT\sim0$, the divisor $T$ is non-trivial torsion, and hence $\mathrm H^0(Y,\OOO_Y(T))=0$. By adjunction (\textit{cf.}~\cite[Corollary~5.25]{KM98}), there is an exact sequence
\[
0\lra\OOO_Y(K_Y)\lra\OOO_Y(T)\lra\OOO_D(K_D)\lra0.
\]
Kawamata--Viehweg vanishing and Serre duality give $\mathrm H^1(Y,\OOO_Y(K_Y))=0$. It follows that $\mathrm H^0(D,\OOO_D(K_D))=0$, which contradicts $K_D\sim0$. Thus $K_Y+D\sim0$.

We have shown that $-K_Y\sim D$. Since $K_Y+D$ is Cartier and $(Y,D)$ is plt, the log discrepancy $a(E,Y,D)$ of every exceptional divisor $E$ over $Y$ is a positive integer. Therefore,
\[
a(E,Y)=a(E,Y,D)+\operatorname{ord}_E(D)\geq1,
\]
so $Y$ has canonical singularities.
 
We claim that there exists a Weil divisor $B$ such that $mB\sim D$. Indeed, let $Y^0$ be a smooth locus of $Y$, and let $X^0$ be its preimage via the map $f$. We denote the corresponding map by $f^0\colon X^0\to Y^0$. Put $D^0=D\cap Y^0$. Then applying \cite[Theorem 1.2]{Wa68}, we see that there exists an element $B^0\in\Pic(Y^0)$ such that $mB^0\sim D^0$. Since $\Pic(Y^0)=\Cl(Y^0)=\Cl(Y)$, we see that $mB\sim D$, where $B$ is the closure of $B^0$ in $Y$. 
\end{proof}

\begin{remark}
\label{rem-bounded-index}
In the above notation, if $Y$ is terminal (this holds for example if the action of $K$ is free outside~$S$), then by \cite{Su04} we have $m\leq 19$. In general, by \cite{JL25} we have $m\leq 66$. 
\end{remark}

\begin{proof}[Proof of Proposition~\ref{thm-second-main-thm}]
This follows from Lemma~\ref{lem-for-thm-2}, Proposition~\ref{prop-K3-boundedness}, and Remark~\ref{rem-bounded-index}.
 \end{proof}

\begin{thebibliography}{AGZV85+++}

\bibitem[AGZV85]{AGV85}
V.\,I.~Arnold, S.\,M.~Gusein-Zade, and A.\,N.~Varchenko,
{\em Singularities of differentiable maps. Vol.~I} (translated from the Russian by I.~Porteous and M.~Reynolds),
Monogr.\ Math., vol.~82, Birkh\"auser, Boston, MA, 1985,
\doi{10.1007/978-1-4612-5154-5}.

\bibitem[Beau07]{Beau07}
A.~Beauville, 
{\em $p$-Elementary subgroups of the Cremona group},
J.~Algebra \textbf{314} (2007), no.~2, 553--564,
\doi{10.1016/j.jalgebra.2005.07.040}.

\bibitem[Bir21]{B21}
C.~Birkar,
{\em Singularities of linear systems and boundedness of Fano varieties},
Ann.\ of Math.~(2) \textbf{193} (2021), no.~2, 347--405,
\doi{10.4007/annals.2021.193.2.1}.

\bibitem[Bir24]{B24}
\bysame, 
{\em On connectedness of non-klt loci of singularities of pairs}, 
J.~Differential Geom.\ \textbf{126} (2024), no.~2, 431--474,
\doi{10.4310/jdg/1712344217}.

\bibitem[Bla07]{Bl07}
J.~Blanc,
{\em Finite Abelian subgroups of the Cremona group of the plane},
C.\ R.\ Acad.\ Sci.\ Paris \textbf{344} (2007), no.~1, 21--26,
\doi{10.1016/j.crma.2006.11.026}.

\bibitem[BH23]{BH23}
S.~Brandhorst and T.~Hofmann, 
{\em Finite subgroups of automorphisms of K3 surfaces}, 
Forum of Math.\ Sigma \textbf{11} (2023), Paper No.~e54,
\doi{10.1017/fms.2023.50}.

\bibitem[DI09]{DI09}
I.\,V.~Dolgachev and V.\,A.~Iskovskikh, 
{\em Finite subgroups of the plane Cremona group},
in: {\em Algebra, arithmetic, and geometry: in honor of Yu.\,I.~Manin. 
Vol.~I}, pp.~443--548, Progr.\ Math., vol.~269, Birkh\"auser Boston,
Boston, MA, 2009,
\doi{10.1007/978-0-8176-4745-2_11}.

\bibitem[dFKX17]{dFKX17}
T.~de Fernex, J.~Koll\'ar, and C.~Xu,
{\em The dual complex of singularities},
in: {\em Higher dimensional algebraic geometry: In honour of Professor Yujiro Kawamata's sixtieth birthday}, pp.~103--129, Adv.\ Stud.\ Pure Math., vol.~74, Math.\ Soc.\ Japan, Tokyo, 2017,
\doi{10.2969/aspm/07410103}.

\bibitem[Ful97]{Fu06} 
W.~Fulton, 
 {\em Young Tableaux: With Application to Representation Theory and Geometry}, 
 London Math.\ Soc.\ Stud.\ Texts, vol.~35, Cambridge Univ.\ Press, Cambridge, 1997, \doi{10.1017/CBO9780511626241}.

\bibitem[Ful00]{Fu00}
\bysame, 
 {\em Eigenvalues, invariant factors, highest weights, and Schubert calculus},
 Bull.\ Amer.\ Math.\ Soc.\ (N.S.) \textbf{37} (2000), no.~3, 209--249,
 \doi{10.1090/S0273-0979-00-00865-X}.

\bibitem[JL25]{JL25}
C.~Jiang and H.~Liu,
{\em A canonical Fano threefold has Fano index $\leq66$}, 
preprint \arXiv{2508.16364v1} (2025).

\bibitem[Kan95]{Kan95}
S.~Kantor,
{\em Theorie der endlichen Gruppen von eindeutigen Transformationen in der Ebene}, Mayer and M\"uller, Berlin, 1895,
\doi{10.1515/crll.1895.114.50}.

\bibitem[Ka69]{Ka69}
I.~Kaplansky,
{\em Infinite Abelian Groups}, revised ed.,
University of Michigan Press, Ann Arbor, MI, 1969.

\bibitem[Kaw88]{Kaw88}
Y.~Kawamata,
 {\em Crepant Blowing-Up of $3$-Dimensional Canonical Singularities and Its Application to Degenerations of Surfaces}, 
 Ann.\ of Math.~(2) \textbf{127} (1988), no.~1, 93--163,
 \doi{10.2307/1971417}.

\bibitem[KM98]{KM98}
  J.~Koll\'{a}r and S.~Mori,
  {\em Birational geometry of algebraic varieties} 
 (with the collaboration of C.\,H.~Clemens and A.~Corti; translated from the 1998 Japanese original), Cambridge Tracts in Math., vol.~134, Cambridge Univ.\ Press, Cambridge, 1998, \doi{10.1017/CBO9780511662560}.

\bibitem[KSB88]{KS-B88}
J.~Koll\'ar and N.\,I.~Shepherd-Barron,
 {\em Threefolds and deformations of surface singularities},
 Invent.\ Math.\ \textbf{91} (1988), no.~2, 299--338,
\doi{10.1007/BF01389370}.
 
 \bibitem[KX16]{KX16}
J. Koll\'{a}r and C.~Xu,
 {\em The dual complex of Calabi--Yau pairs}, 
 Invent.\ Math.\ \textbf{205} (2016), no.~3, 527--557,
 \doi{10.1007/s00222-015-0640-6}.

\bibitem[KZ26]{KZh24}
J.~Koll\'{a}r and Z.~Zhuang.
{\em Essential dimension of isogenies},
Pure Appl.\ Math.~Q.\ \textbf{22} (2026), no.~2, 675--690,
\doi{10.4310/pamq.250306163938}.


\bibitem[Kuz20]{Kuz20}
A.\,A.~Kuznetsova, 
 {\em Finite $3$-Subgroups in the Cremona Group of Rank $3$},
 Math.\ Notes \textbf{108} (2020), no.~5-6, 697--715,
 \doi{10.1134/S0001434620110085}.

\bibitem[Log22]{Lo22}
K.~Loginov, 
 {\em A note on 3-subgroups in the space Cremona group}, 
 Comm.\ Algebra \textbf{50} (2022), no.~9, 3704--3714,
 \doi{10.1080/00927872.2022.2042545}.

\bibitem[Mor85]{Mo85}
S.~Mori, 
 {\em On $3$-dimensional terminal singularities}, 
 Nagoya Math.~J.\ \textbf{98} (1985), 43--66,
 \doi{10.1017/S0027763000021358}.

\bibitem[Pop14]{Po14}
V.\,L.~Popov,
 {\em Jordan groups and automorphism groups of algebraic varieties}, 
 in: {\em Automorphisms in birational and affine geometry}, pp.~185--213, Springer Proc.\ Math.\ Stat., vol.~79, Springer, Cham, 2014,
 \doi{10.1007/978-3-319-05681-4_11}.
  
\bibitem[Pro09]{Pr09}
Yu.~Prokhorov, 
 {\em Simple finite subgroups of the Cremona group of rank $3$}, 
 J.~Algebraic Geom.\ \textbf{21} (2009), no.~3, 563--600,
 \doi{10.1090/S1056-3911-2011-00586-9}.

\bibitem[Pro11]{Pr11}
\bysame, 
{\em {$p$}-elementary subgroups of the {C}remona group of rank $3$}, 
in: {\em Classification of algebraic varieties}, pp.~327--338,
EMS Ser.\ Congr. Rep., Eur.\ Math.\ Soc., Z\"urich, 2011,
\doi{10.4171/007-1/16}.

\bibitem[Pro14]{Pr14}
\bysame, 
 {\em $2$-elementary subgroups of the space {C}remona group},
 in: {\em Automorphisms in birational and affine geometry}, pp.~215--229, 
 Springer Proc.\ Math.\ Stat., vol.~79, Springer, Cham, 2014,
 \doi{10.1007/978-3-319-05681-4_12}.

\bibitem[PS16]{PS16}
Yu.~Prokhorov and C.~Shramov, 
 {\em Jordan property for Cremona groups},
 Amer.\ J.~Math.\ \textbf{138} (2016), no.~2, 403--418,
 \doi{10.1353/ajm.2016.0017}.

\bibitem[PS18]{PS17}
\bysame, 
 {\em $p$-subgroups in the space Cremona group},
 Math.\ Nachr.\ \textbf{291} (2018), no.~8-9, 1374--1389, 
 \doi{10.1002/mana.201700030}.

\bibitem[Rei87]{Re85}
M.~Reid, 
{\em Young person's guide to canonical singularities},
in: \emph{Algebraic  Geometry, {B}owdoin, 1985} ({B}runswick, {M}aine, 1985), pp.~345--414, Proc.\ Sympos.\ Pure  Math., vol.~46, Part~1, Amer.\ Math.\ Soc., Providence, RI, 1987, \doi{10.1090/pspum/046.1/927963}.

\bibitem[Suz04]{Su04}
K.~Suzuki, 
 {\em On Fano indices of $\mathbb Q$-Fano 3-folds},
 Manuscripta Math.\ \textbf{114} (2004), no.~2, 229--246,
 \doi{10.1007/s00229-004-0442-4}. 

\bibitem[Wav68]{Wa68}
J.\,J.~Wavrik.
 {\em Deformations of branched coverings of complex manifolds}.
 Amer.~J.\ Math.\ \textbf{90} (1968), 926--960,
 \doi{10.2307/2373491}.

\bibitem[Wim96]{Wim96}
A.~Wiman, 
 {\em Zur Theorie der endlichen Gruppen von birationalen Transformationen in der Ebene}, 
 Math.\ Ann.\ \textbf{48} (1896), no.~1-2, 195--240,
 \doi{10.1007/BF01446342}.

\bibitem[Xu20]{Xu20}
J.~Xu, 
 {\em A remark on the rank of finite $p$-groups of birational
automorphisms}, 
 C.\ R.\ Math.\ Acad.\ Sci.\ Paris \textbf{358} (2020), no.~7, 827--829,
 \doi{10.5802/crmath.93}. 

\bibitem[Zim12]{Zim12}
B.\,P.~Zimmermann, 
 {\em On finite groups acting on spheres and finite subgroups of orthogonal groups}, 
 Sib.\ \`Elektron.\ Mat.\ Izv.\ \textbf{9} (2012), 1--12.

\end{thebibliography}
\end{document}
